\section{Hardness on graphs of bounded degree}
\label{sec:bounded-degree-hardness}
\label{sec:expander-sparsification}
\label{sec:bounded-degree-treewidth}

\subsection{The degree-reduction approach}

Let $H$ be one of the dense graphs constructed above, with
$N=|V(H)|$. We describe its edges as a union of specified
\emph{factors} and an edge set $E_0$ that we retain. A clique factor
on $U$ contributes every pair in $\binom U2$. A \emph{biclique}
factor with disjoint sides $A,B$ contributes every edge $uv$ with
$u\in A$ and $v\in B$. Factors may overlap. A factor specifies
edges to be included, rather than an induced subgraph: other factors
or $E_0$ may add further edges on the same vertices.

We replace each factor by a graph of bounded degree on the same
vertices. Write $J$ for the union of the replacements
and $E_0$, so $J$ is a spanning subgraph of $H$. The argument has
three steps.
\begin{enumerate}
 \item \emph{Preserve the YES bound by deleting edges.}
   Every elimination forest or tree decomposition of $H$ is also
   valid for $J$. Thus none of the three optimum values increases.
 \item \emph{Recover a dense solution from a sparse one.}
   Each replacement has an edge between every two subsets that meet
   its chosen threshold. For a clique factor the subsets are disjoint,
   and for a biclique they lie on opposite sides.
   The completion lemmas show that, after deleting a small set $D$
   of exceptional vertices, an arbitrary forest or decomposition
   of $J$ also represents every edge of $H-D$. For forests,
   restore $D$ as a chain above all existing roots. For tree
   decompositions, add $D$ to every bag. The height or width
   increases by at most $|D|$, and the total depth by at most
   $N|D|$. Thus a sparse solution of value $z$ gives a dense
   solution of height or width at most $z+|D|$, or total depth at
   most $z+N|D|$.
 \item \emph{Keep the total loss below the hardness gap.}
   Choose each factor's threshold as a small fraction of its
   number of vertices. If each vertex belongs to at most a fixed
   number of factors, the sum of all factor sizes is $O(N)$.
   Hence $|D|\le\varepsilon N$ for any prescribed fixed
   $\varepsilon>0$, by choosing those fractions sufficiently small.
   The resulting losses are at most $\varepsilon N$ for treedepth
   and treewidth, and $\varepsilon N^2$ for total depth. The dense
   gaps have these scales. If their leading YES and NO coefficients
   are $A$ and $B$ with $B>rA$, choose $\varepsilon$ so the loss
   coefficient is below $B-rA$.
\end{enumerate}

For the maximum degree to be bounded, we also need bounded degree
in $E_0$ and a bounded ratio between the two side sizes of every
biclique. The expander construction below then gives a constant
degree contribution from each factor. Bounded factor membership
bounds their sum at every vertex. These constants may depend on
the desired approximation ratio, but not on $N$. For treewidth,
each biclique side must additionally lie in a bounded number of
clique factors. These supply the bags used in its completion proof.
The UGC applications require a bounded-degree independent-set source so
that their retained source edges satisfy the condition on $E_0$.

\subsection{Completing sparse factors}

Two disjoint vertex sets are \emph{anticomplete} if no edge joins them.
A \emph{$(t,t)$ hole} is a pair of anticomplete sets of $t$ vertices
each. For a bipartite replacement on $A\mathbin{\dot\cup}B$, we
consider only holes with one set in each side. Thus it has no
$(t,t)$ hole if some edge joins $S$ to $T$ whenever
$S\subseteq A$, $T\subseteq B$, and $|S|=|T|=t$.
For a clique replacement on $U$, we consider all disjoint
$t$-subsets $S,T\subseteq U$.
Thus the condition is inherited by deleting vertices. Recall that
forest restriction preserves the ancestor order between surviving
vertices and never increases their depths.

The following lemmas identify a small exceptional set whose deletion
allows a sparse factor to be completed. For a clique, the surviving
vertices must lie on one ancestor path or in one bag. For a biclique,
it is enough that every surviving cross-pair is comparable or appears
together in a bag.

\subsubsection{Clique factors}

\begin{lemma}[Clique completion in an elimination forest]\label{lem:sparse-clique-completion}
Let $F$ be an elimination forest of a graph $J$. Suppose $J[U]$ has
no two disjoint anticomplete sets of size $t\ge1$. An ancestor path
of $F$ contains all but at most $2(t-1)$ vertices of $U$.
\end{lemma}

\begin{proof}
If $|U|<t$, the empty path suffices. Otherwise traverse $F$ in
depth-first order, traversing its rooted trees successively.
Introduce a vertex before visiting its descendants, and forget it
after traversing its entire subtree. The introduced but unforgotten
vertices form an ancestor path at every time. Every forgotten vertex
is incomparable with every unintroduced vertex: all its ancestors
were introduced earlier, and its whole subtree has already been
traversed. Thus the two groups are anticomplete in $J$.

Immediately after the $t$th vertex of $U$ is forgotten, at most
$t-1$ vertices of $U$ are unintroduced: otherwise they and the
$t$ forgotten vertices would form a forbidden anticomplete pair.
Immediately before that forgetting operation, $t-1$
vertices of $U$ had been forgotten, and the unintroduced set was
the same. All other vertices of $U$ lay on the current ancestor
path. That path therefore contains at least $|U|-2(t-1)$ vertices
of $U$.
\end{proof}

\begin{lemma}[Clique completion in a tree decomposition]
\label{lem:tw-clique-completion}
Suppose $J[U]$ has no two disjoint anticomplete sets of size
$t\ge1$. Every tree decomposition of $J$ has a bag containing
all but at most $3(t-1)$ vertices of $U$.
\end{lemma}

\begin{proof}
Let $T$ be the decomposition tree. For each graph vertex $v$, let
$T_v$ be the nonempty connected subtree of nodes whose bags contain
$v$.
If $|U|<t$, any bag suffices. Otherwise consider an edge of the
decomposition tree. Its deletion leaves two sides. Call a side
heavy if it contains the entire subtree $T_v$ for at least $t$
vertices $v\in U$. Both sides cannot be heavy: their corresponding
vertex sets would be anticomplete. Orient the edge toward its
heavy side if one exists, and arbitrarily if neither is heavy.
Following outgoing edges in a finite tree eventually reaches a
node with no outgoing edge. Call such a node a \emph{sink}, and let
$u$ be one.

Each component of $T-u$ contains the entire subtree of at most
$t-1$ vertices of $U$, by the sink property. Every vertex of $U$
absent from the bag at $u$ belongs wholly to one of these components.
For $t=1$ there are no such vertices. For $t\ge2$, suppose their
total were at least $3t-2$. Collect components until the number of
vertices of $U$ assigned to them first reaches $t$. This count is at most
$2t-2$, since each component contributes at most $t-1$. At least
$t$ vertices remain in the other components. The two groups are
anticomplete, a contradiction. Thus at most $3t-3$ vertices of
$U$ are absent from the bag at $u$.
\end{proof}

\subsubsection{Biclique factors}

We use K\H onig's theorem~\cite[Theorem~2.1.1]{Diestel2017}: in a
bipartite graph, the maximum matching size equals the minimum
vertex-cover size.

\begin{lemma}[Biclique completion in an elimination forest]\label{lem:sparse-biclique-completion}
Let $F$ be an elimination forest of a graph $J$. Suppose $A,B$ are
disjoint and the bipartite graph of $J$ between them has no $(t,t)$
hole, where $t\ge1$. There is a set $D\subseteq A\cup B$ with
$|D|\le4(t-1)$ such that every pair in
$(A\setminus D)\times(B\setminus D)$ is comparable in $F$.
\end{lemma}

\begin{proof}
A depth-first traversal introduces a vertex, traverses each of its
child subtrees completely, and then forgets it. Assign a distinct
time to each introduction and forgetting event, traversing different
rooted trees consecutively. Let $t_{\mathrm{in}}(v)$ and $t_{\mathrm{out}}(v)$ be the
introduction and forgetting times of $v$. Its descendants are the
vertices visited during $[t_{\mathrm{in}}(v),t_{\mathrm{out}}(v)]$.
Two vertices are incomparable if and only if these intervals
are disjoint. Let $M$ be the bipartite graph on $A,B$ whose edges
are the incomparable pairs. Partition its edges according to
$t_{\mathrm{out}}(u)<t_{\mathrm{in}}(v)$ or $t_{\mathrm{out}}(v)<t_{\mathrm{in}}(u)$.
If $M$ has a matching of size at least $4t-3$, one type contains
at least $2t-1$ edges of that matching. For the first type, write
these edges as $u_iv_i$, $1\le i\le m_0$, with $m_0\ge2t-1$ and
$t_{\mathrm{out}}(u_1)\le\cdots\le t_{\mathrm{out}}(u_{m_0})$. If $i\le t$ and $j\ge m_0-t+1$,
then $i\le j$, and hence $t_{\mathrm{out}}(u_i)\le t_{\mathrm{out}}(u_j)<t_{\mathrm{in}}(v_j)$.
The first $t$ vertices $u_i$ and the last $t$ vertices $v_j$ are
therefore pairwise incomparable across the two sets. They are
anticomplete in $J$, a contradiction. The second type is symmetric.
It follows that the maximum matching size of $M$ is at most $4t-4$.
K\H onig's theorem supplies a vertex cover $D$ of $M$ of that size.
All surviving cross-pairs are comparable.
\end{proof}

For treewidth, we will first complete the clique factors and then
use their bags to anchor the sides of each biclique. Restricting
vertex subtrees to the path between two such bags reduces the
argument to prefixes and suffixes, using the subtree projection
in~\cite[proof of Lemma~5]{BonnetTreewidth}. Here the no-hole
condition bounds a matching of uncovered pairs. A vertex cover of
those pairs gives the required deletion set.

When vertices are deleted from a graph, we restrict its tree
decomposition by removing them from every bag.

\begin{lemma}[Biclique completion between two bags]
\label{lem:tw-biclique-completion}
Fix a tree decomposition of $J$. Suppose disjoint sets $A,B$
are contained in bags at nodes $x,y$, respectively, and the
bipartite graph of $J$ between $A$ and $B$ has no $(t,t)$ hole.
One can delete at most $2(t-1)$ vertices of $A\cup B$ so that
every edge of the complete bipartite graph between the surviving
sets is already covered by a bag of the restricted decomposition.
\end{lemma}

\begin{proof}
For each vertex $v$, let $T_v$ be the subtree of decomposition
nodes whose bags contain $v$.
If $x=y$, no deletion is needed. Otherwise let $P$ be their
connecting path. For $u\in A$, the intersection $T_u\cap P$
is a nonempty prefix of $P$, starting at $x$. For $v\in B$,
$T_v\cap P$ is a suffix ending at $y$.
The subtrees $T_u,T_v$ intersect if and only if these path
intervals intersect. For the nontrivial direction, if both
contain some node $z$, they both contain the median of $x,y,z$,
which lies on $P$.

Let $M$ be the bipartite graph on $A,B$ in which $uv$ is an edge
if and only if $T_u$ and $T_v$ are disjoint.
If it has a matching of size $m_0\ge2t-1$, write its edges
$u_iv_i$ in increasing order of the endpoint of the prefix for
$u_i$. The endpoint of each $u_i$-prefix precedes the start
of its matched $v_i$-suffix. Consequently every one of the first
$t$ prefixes is disjoint from every one of the last $t$ matched
suffixes: if $i\le t$ and $j\ge m_0-t+1$, then $i\le j$.
Their corresponding $t$-subsets of $A$ and $B$ are anticomplete
in $J$, contradicting the hypothesis. Thus the maximum matching size
of $M$ is at most $2t-2$.

K\H onig's theorem gives a vertex cover of $M$ of size at most
$2t-2$. Delete these vertices. Every surviving cross-pair has
intersecting bag subtrees, so its completed edge is covered.
\end{proof}

\subsubsection{Comparison bounds}

We now sum the exceptions over all factors. If one vertex is deleted
for several factors, it appears only once in the union of the deletion
sets. The treedepth and total-depth bounds apply to
arbitrary factors. The treewidth bound additionally uses a containing clique
for each side of a biclique, providing the two bags required above.

\begin{lemma}[Treedepth comparison]\label{lem:factor-sparsification}
Suppose $H$ is the union of an edge set $E_0$, clique factors
on $U_f$ indexed by $f\in\mathcal C$, and complete bipartite
factors $A_f\times B_f$ indexed by $f\in\mathcal B$.
For each factor choose a positive integer $t_f$ and replace it by
a subgraph with no $(t_f,t_f)$ hole. For a clique factor, this means
that every two disjoint $t_f$-subsets have an edge between them. For a
biclique factor, it means that every two $t_f$-subsets on opposite
sides have an edge between them. Retain $E_0$ and let $J$ be the resulting graph on
$V(H)$. Put $M=2\sum_{f\in\mathcal C}(t_f-1)
   +4\sum_{f\in\mathcal B}(t_f-1)$.
Every elimination forest $F$ of $J$ can be completed to an
elimination forest $\widehat F$ of $H$ with
$h(\widehat F)\le h(F)+M$.
In particular, $\td(J)\le\td(H)\le\td(J)+M$.
Factors may overlap in vertices or edges.
\end{lemma}

\begin{proof}
Fix any elimination forest $F$ of $J$.
For each clique factor delete the vertices
outside the path given by Lemma~\ref{lem:sparse-clique-completion}.
For each biclique factor use the deletion set in
Lemma~\ref{lem:sparse-biclique-completion}. Let $D$ be their union.
Each vertex of this union is counted at least once in the sum
defining $M$, so $|D|\le M$, even when factors overlap.
In the restriction of $F$ to $V(H)\setminus D$,
every surviving edge of every completed factor has comparable
endpoints. Every retained edge also has comparable endpoints.
Thus the restriction is an elimination forest of $H-D$, and its
height is at most $h(F)$. By Lemma~\ref{lem:restore-vertices},
restoring $D$ gives an elimination forest $\widehat F$ of $H$ with
$h(\widehat F)\le h(F)+|D|\le h(F)+M$.
Choosing $F$ of minimum height proves
$\td(H)\le\td(J)+M$. The other comparison follows from
$J\subseteq H$.
\end{proof}

\Needspace{11\baselineskip}
\begin{lemma}[Total-depth comparison]
\label{lem:total-depth-comparison}
Let $H,J$, the factor thresholds, and $M$ be as in
Lemma~\ref{lem:factor-sparsification}, and put $N=|V(H)|$.
Every elimination forest $F$ of $J$ can be completed to an
elimination forest $\widehat F$ of $H$ with
$\totd(\widehat F)\le\totd(F)+NM$.
In particular, $\totd(J)\le \totd(H)\le \totd(J)+NM$.
\end{lemma}

\begin{proof}
Fix any elimination forest $F$ of $J$. Use the exceptional set
$D$ from the proof of Lemma~\ref{lem:factor-sparsification} for
this forest, so $|D|\le M$. That proof shows that the
restriction $F'$ of $F$ to $V(H)\setminus D$ is an elimination
forest of $H-D$.

Deleting vertices and restricting the forest can only decrease the
depth of a surviving vertex, so $\totd(F')\le\totd(F)$. Restore
the vertices of $D$ as in Lemma~\ref{lem:restore-vertices}, placing
them one at a time above all existing roots. Each restoration
increases the depths of the existing vertices by one and gives
the new root depth one. Its total cost is therefore the number of
vertices then present, at most $N$. After $|D|$ restorations, the
resulting elimination forest $\widehat F$ satisfies
$\totd(\widehat F)\le\totd(F)+N|D|\le\totd(F)+NM$.
Choosing $F$ of minimum total depth proves the upper
comparison. The lower comparison follows from $J\subseteq H$.
\end{proof}

\begin{lemma}[Treewidth comparison]
\label{lem:tw-factor-sparsification}
Let $H$ be the union of retained edges $E_0$, clique factors
on $U_f$ indexed by $f\in\mathcal C$, and complete bipartite
factors $A_f\times B_f$ indexed by $f\in\mathcal B$.
Suppose each side of every biclique factor is contained in
$U_f$ for some $f\in\mathcal C$. For each factor choose a
positive integer threshold $t_f$. Replace each factor by a
subgraph with no hole at its chosen threshold, retain $E_0$,
and call the resulting spanning graph $J$.
Put $M=3\sum_{f\in\mathcal C}(t_f-1)
   +2\sum_{f\in\mathcal B}(t_f-1)$.
Every tree decomposition of $J$ can be completed to one of $H$
whose width increases by at most $M$. In particular,
$\operatorname{tw}(J)\le\operatorname{tw}(H)
 \le\operatorname{tw}(J)+M$.
Clique and biclique factors may overlap in vertices or edges.
\end{lemma}

\begin{proof}
Fix any tree decomposition of $J$. For each clique
factor apply Lemma~\ref{lem:tw-clique-completion}, and delete
the exceptions outside its chosen bag. In the restricted
decomposition each surviving set $U_f$ is contained in
its chosen bag, so its clique can be completed.

Each surviving biclique factor still has no $(t_f,t_f)$ hole,
and its sides are contained in the chosen bags of their containing
clique factors.
Apply Lemma~\ref{lem:tw-biclique-completion} to choose at most
$2(t_f-1)$ further exceptions for each factor. Let $D$ be the
union of all exception sets. Deleting other factors' exceptions
does not remove a bag intersection between surviving vertices.
Thus the restricted decomposition covers every edge of $H-D$,
including all completed factors and all retained edges.
Add every vertex of $D$ to every bag. The resulting bags form a
decomposition of $H$: every edge with an endpoint in $D$ is covered,
and the bags containing each vertex remain connected. Its width exceeds
that of the supplied decomposition by at most $|D|\le M$.
Choosing a minimum-width decomposition proves the upper comparison.
The lower comparison follows from $J\subseteq H$.
\end{proof}

\subsection{Sparse replacements at arbitrary sizes}

We next construct bounded-degree graphs satisfying the no-hole
condition. We first introduce the expander facts used in the
construction, then map an expander onto each prescribed factor.

\subsubsection{Expanders, walks, and the mixing lemma}

The replacement graphs have a fixed degree. We measure how evenly
their edges join vertex sets through the eigenvalues of the adjacency
matrix. The following definitions and mixing lemma make
this condition precise. A \emph{step} traverses one edge in a chosen
direction. A loop step stays at the same vertex.
For an intermediate undirected multigraph
$\mathcal E$ on $M$ vertices, let $A$ be its symmetric adjacency
matrix. The entry $A_{uv}$ counts the available steps from $u$ to
$v$, including loop steps when $u=v$. Parallel edges give separate
steps. In this subsection, $d$-regular means that every row of $A$
sums to $d$, so there are $d$ possible next steps at every vertex.
The all-ones vector $\mathbf1$ is then an eigenvector with eigenvalue
$d$. A \emph{spectral expander with eigenvalue bound $\lambda<d$}
means that every eigenvalue on the subspace
$\{x:\sum_vx_v=0\}$ has absolute value at most $\lambda$.
For a bounded-degree expander family, both $d$ and an upper bound
smaller than one on $\lambda/d$ are fixed independently of $M$.

For $S,T\subseteq V(\mathcal E)$, put
$e_{\mathcal E}(S,T)=\sum_{u\in S,v\in T}A_{uv}$. This counts
ordered steps with multiplicity. The sets need not be disjoint.
The following lemma states the distribution property we need.

\begin{lemma}[Expander mixing lemma {\cite[Lemma~2.5]{HLW2006}}]
\label{lem:expander-mixing}
For every $d$-regular $\mathcal E$ with eigenvalue bound $\lambda$,
and all $S,T\subseteq V(\mathcal E)$,
\[
 \left|e_{\mathcal E}(S,T)-\frac{d|S||T|}{M}\right|
 \le\lambda\sqrt{|S||T|}.
\]
\end{lemma}

\begin{proof}
The spectral proof of the cited lemma also applies to our step
matrix, including loops and parallel edges.
Let $\mathbf1_S$ be the vector equal to one on $S$ and zero
elsewhere, and put $x_S=\mathbf1_S-(|S|/M)\mathbf1$. Define $x_T$
similarly. These two vectors have coordinate sum zero. Since
$A\mathbf1=d\mathbf1$,
$e_{\mathcal E}(S,T)-d|S||T|/M=x_S^{\mathsf T}Ax_T$.
The eigenvalue bound and Cauchy--Schwarz give absolute value at most
$\lambda\|x_S\|_2\|x_T\|_2$, where $\|\cdot\|_2$ is the Euclidean
norm. Since $\|x_S\|_2^2=|S|(1-|S|/M)\le|S|$, and similarly for
$T$, this is at most $\lambda\sqrt{|S||T|}$.
\end{proof}

A \emph{walk of length $k$} consists of a starting vertex followed
by $k$ consecutive edge steps. Vertices and edges may be repeated,
and different choices among parallel edges count as different
walks. The \emph{$k$th walk power} of $\mathcal E$ has the same
vertices, with one step from $u$ to $v$ for each length-$k$ walk
from $u$ to $v$ in $\mathcal E$. Its adjacency matrix is $A^k$.
It is $d^k$-regular, and its eigenvalues on the zero-sum subspace
have absolute value at most $\lambda^k$. Thus taking a fixed walk
power makes the ratio $(\lambda/d)^k$ as small as needed while
keeping the degree independent of $M$.

The Margulis--Gabber--Galil construction gives, for every integer
$s\ge1$, a deterministically constructible eight-regular multigraph
on $\mathbb Z_s^2$, where $\mathbb Z_s$ denotes the integers modulo
$s$, with eigenvalue bound $5\sqrt2$
\cite[Construction~8.1 and Theorem~8.2]{HLW2006}. Its intermediate
loops and parallel edges are allowed under the convention above.
Put $\lambda_0=5\sqrt2/8<1$. Its $k$th walk power has
$M=s^2$ vertices, degree $d=8^k$, and eigenvalue bound
$d\lambda_0^k$. Lemma~\ref{lem:expander-mixing} therefore gives
\begin{equation}\label{eq:factor-expander-mixing}
 e_{\mathcal E}(S,T)\ge d|S||T|/M-d\lambda_0^k\sqrt{|S||T|}
 \quad\text{for all }S,T,
\end{equation}
where $\mathcal E$ now denotes that walk power.
In particular, it joins any two sets of size at least $t$ whenever
$t>\lambda_0^kM$.

\subsubsection{Mapping to the prescribed vertex sets}

\begin{lemma}[Sparse replacements on prescribed vertex sets]\label{lem:factor-expanders}
For integers $n_A,n_B\ge1$, set
$M=\lceil\sqrt{\max(n_A,n_B)}\rceil^2$.
For each fixed integer $k\ge1$ there is a deterministic polynomial-time
construction of a simple bipartite graph on prescribed disjoint
sets of sizes $n_A,n_B$ with maximum degree at most
$8^k\max\{\lceil M/n_A\rceil,\lceil M/n_B\rceil\}$,
having no $(t,t)$ hole whenever $t>\lambda_0^kM$.
For one set of size $n_A$, the same construction gives a simple
graph of maximum degree at most $8^k\lceil M/n_A\rceil$,
where $M=\lceil\sqrt{n_A}\rceil^2$, with no disjoint anticomplete
$t$-subsets whenever $t>\lambda_0^kM$.
Neither construction adds vertices.
\end{lemma}

\begin{proof}
Let $A,B$ be the prescribed sides, and let $\mathcal E$ be the
walk power on $M$ vertices just described. Map its vertices onto
$A$ by $\pi_A$ so that every vertex of $A$ receives either
$\lfloor M/n_A\rfloor$ or $\lceil M/n_A\rceil$ intermediate vertices.
Choose a map $\pi_B$ onto $B$ in the same way, using the corresponding
side size $n_B$. Such maps exist because $M\ge\max\{n_A,n_B\}$.
Put an edge $uv$, with $u\in A$ and $v\in B$, in the output if
and only if some step $xy$ of $\mathcal E$ satisfies
$\pi_A(x)=u$ and $\pi_B(y)=v$.
Repeated occurrences give a single output edge.
Each intermediate vertex has $8^k$ steps leaving it and, by
symmetry, $8^k$ steps entering it. Consequently an output vertex
receives at most $8^k$ incident edges per intermediate vertex
mapped to it, proving the degree bound.

For subsets $S_A\subseteq A$ and $T_B\subseteq B$ of size $t$,
let $S=\pi_A^{-1}(S_A)$ and $T=\pi_B^{-1}(T_B)$ be the sets of
intermediate vertices mapped into them. Both maps are onto, so
$|S|,|T|\ge t$. Equation~\eqref{eq:factor-expander-mixing}
gives a step between them when $t>\lambda_0^kM$, since
\[
 |S||T|/M-\lambda_0^k\sqrt{|S||T|}
 =\sqrt{|S||T|}\big(\sqrt{|S||T|}/M-\lambda_0^k\big)>0.
\]
That step produces an output edge between $S_A$ and $T_B$.
The sets $S,T$ may overlap. The mixing inequality still applies. An
intermediate loop also gives a legitimate edge between the
two disjoint output sides.

For the one-set construction, map the intermediate vertices onto
the prescribed set by $\pi$, with at most $\lceil M/n_A\rceil$
vertices mapped to any output vertex. Include an edge $uv$ for
distinct output vertices $u,v$ if and only if a step of $\mathcal E$
has its endpoints mapped to $u,v$.
Discard loops and repeated edges. Intermediate vertices mapped into
disjoint output subsets form disjoint sets, so an edge between them does not
disappear as a loop. In either construction, $M$ is at
most four times the largest prescribed set size, so the construction
has polynomial size for fixed $k$.
\end{proof}

\subsection{A bounded-degree independent-set source under UGC}
\label{sec:bounded-degree-ugc-source}

We derive a bounded-degree version of the independent-set promise
used in Section~\ref{sec:ugc}: independent-set density close to
$1/2$ in the YES case and arbitrarily small constant density in
the NO case. We use the weighted density construction of Austrin,
Khot, and Safra~\cite[Theorem~3.1]{AKS2011}. Their bounded-degree
theorem uses randomized edge sampling and removal of high-degree
vertices~\cite[Theorem~4.1 and Section~4, Steps~2--3]{AKS2011}.
We instead use deterministic spectral sparsification
\cite[Theorem~1.1]{BSS2012} and discard edges below a fixed weight
threshold. The resulting degree bound depends only on the fixed source
gap parameters.

The derivation has three steps. We first remove loops while retaining
the weighted density promise. We then replace each vertex by
independent copies in proportion to its weight, obtaining a
uniform bound on weighted degrees. Finally, a sparse reweighted
graph preserves positive edge
weight inside every vertex set above a fixed size threshold. Discarding
edges turns the weighted-degree bound into an ordinary degree bound.

\subsubsection{The weighted density source and removal of loops}

Let $\Psi$ be the cumulative distribution function of a standard
normal variable, with mean zero and variance one. The value $\Psi^{-1}(s)$
is its $s$-quantile. For $-1<\xi<1$, take jointly normal variables
$X,Y$ of mean zero, variance one, and covariance $\xi$, and define
$\Gamma_\xi(s)=\Pr(X\le\Psi^{-1}(s),\,Y\le\Psi^{-1}(s))$.
Set $\Gamma_\xi(0)=0$ and $\Gamma_\xi(1)=1$.
This function is positive for $0<s<1$, nondecreasing, and convex.
Indeed, differentiating the joint distribution gives
$\Gamma_\xi'(s)=2\Psi(\sqrt{(1-\xi)/(1+\xi)}\,\Psi^{-1}(s))$,
which is nonnegative and increasing. These are the properties used
in the density analysis of~\cite[Theorem~3.1 and equation~(3.2)]{AKS2011}.

A bipartite \emph{Unique Games instance} has left vertices $U$,
right vertices $V$, and a finite label set, called the
\emph{alphabet}, $[\ell]$. Each edge $uv$ carries a permutation
$\pi$ of $[\ell]$ and requires the label at $v$ to be $\pi$ applied
to the label at $u$. A \emph{labeling} assigns one label to every
vertex. Its value is the fraction of constraints it satisfies,
counting parallel constraints separately. The instance's optimum
is the largest such fraction.
We use the regular bipartite formulation in
\cite[Conjecture~2.4]{AKS2011}, which is equivalent to the ordinary
UGC by Khot and Regev~\cite{KhotRegev}. In its YES case, a
$(1-\gamma)$-fraction of the right vertices admit labels satisfying
every incident constraint. Its NO case has optimum at most
$\gamma$. The underlying bipartite multigraph is regular in
both cases: every left vertex has one common degree, and every
right vertex has another common degree. This final condition
in~\cite[Conjecture~2.4]{AKS2011} supplies the weighted-degree
identity used below in both cases.

Fix a rational $p\in(0,1/2)$ and put $\xi=-p/(1-p)$.
Theorem~3.1 of~\cite{AKS2011} gives a deterministic polynomial-time
reduction, for fixed alphabet size, to a vertex- and edge-weighted
graph $G$ with both total weights equal to one. If $a_v$ is the
vertex weight, write $a(S)=\sum_{v\in S}a_v$. The weighted degree
of $v$, the sum of its incident edge weights, is $2a_v$. In the YES case,
$G$ has an independent set of weight at least $p-\gamma$. In the
NO case, $e_G(S)\ge\Gamma_\xi(a(S))-\varepsilon_{\mathrm{dens}}$ for every
$S\subseteq V(G)$. Here $\varepsilon_{\mathrm{dens}}>0$ is an arbitrarily small fixed analytic error
and the Unique Games error $\gamma$ is chosen sufficiently small.
Here $e_G(S)$ is the internal edge weight. We allow loops in this
intermediate graph, counting a loop once in $e_G(S)$ and twice in
weighted degree.

We remove loops before replacing vertices by independent copies so that the NO-case
argument can compare internal edge weights through independent
choices of distinct original vertices. The following alphabet
padding makes the total loop weight arbitrarily small.

\begin{lemma}[Removing loops by alphabet padding]\label{lem:aks-loop-removal}
For every fixed $\varepsilon_{\mathrm{loop}}>0$, a fixed alphabet padding followed by the
AKS construction and deletion of all loops gives a graph $G^-$
with vertex weights $a_v$ of total one, total edge weight at most
one, and weighted degrees at most $2a_v$.
Its YES independent-set weight is at least $p-\gamma$, and in the
NO case
\begin{equation}\label{eq:aks-loop-free-density}
 e_{G^-}(S)\ge\Gamma_\xi(a(S))-\varepsilon_{\mathrm{dens}}-\varepsilon_{\mathrm{loop}}
 \qquad\text{for every }S.
\end{equation}
The reduction is deterministic and polynomial for fixed parameters.
\end{lemma}

\begin{proof}
The edge distribution in~\cite[proof of Theorem~3.1]{AKS2011} first chooses two
constraints incident with the same left vertex, then samples
independent coordinate pairs $(x_j,y_j)$ with probabilities
$\Pr(0,0)=1-2p$, $\Pr(0,1)=\Pr(1,0)=p$, and $\Pr(1,1)=0$.
The endpoint strings are then permuted by the two constraints.
To control possible loops even when the Unique Games instance has
parallel constraints with different permutations, we first pad its
alphabet from $[\ell]$ to $[\ell]\times[t]$, replacing every
constraint $\pi$ by $(i,j)\mapsto(\pi(i),j)$.
Projecting any satisfying padded label pair to its first coordinate
shows that padding cannot increase the NO optimum. A YES labeling
extends by giving every labeled vertex the same second coordinate.
Regularity and the YES fraction are unchanged.

The coordinates $[\ell]\times\{j\}$, for a fixed $j\in[t]$,
form one coordinate block. Put $c=p^2+(1-p)^2<1$. Conditional on any choice of the two
constraints, a loop requires equality of the two permuted strings
in every one of the $t$ coordinate blocks. In one block, consider
one output coordinate. If the two permutations select the same
input coordinate, equality has probability $1-2p$. If they select
different input coordinates, the two bits are independent with
marginal probability $p$, so equality has probability $c$.
In either case it is at most $c$. Different blocks are independent.
Consequently the total loop weight is at most $c^t$.
Choose a fixed $t$ with $c^t<\varepsilon_{\mathrm{loop}}$, and delete all loops.
This loses at most $\varepsilon_{\mathrm{loop}}$ internal edge weight in any set,
preserves the YES independent set, and does not increase weighted
degrees. Alphabet padding preserves the source promise and
regularity, so the construction above has the stated bounds.
All operations are deterministic, and $t$ is fixed.
\end{proof}

\subsubsection{Vertex unweighting and uniform degree bounds}

Write the padded alphabet size as $R$ and the fixed bias as
$p=a_0/a_1$, with positive integers $a_0<a_1$. For a binary
string $x\in\{0,1\}^R$, let $|x|$ be its
number of coordinates equal to one. The vertex weights in the AKS
construction, indexed by right vertices $v\in V$ and strings $x$,
have the form
\[
 a_{(v,x)}=
 \frac{a_0^{|x|}(a_1-a_0)^{R-|x|}}{|V|a_1^R},
 \qquad v\in V,\quad x\in\{0,1\}^R.
\]
Thus $N=|V|a_1^R$ is polynomial in the Unique Games instance, since
$a_1$ and $R$ are fixed. Replace each original vertex $u$ by
$m_u=Na_u$ independent copies, called \emph{clones}, labeled
$(u,1),\ldots,(u,m_u)$. These copies form the \emph{clone block} of $u$.
For every edge $uv$ of weight $w_{uv}$ in $G^-$, join each clone
of $u$ to each clone of $v$ by an edge of weight
$w_{uv}/(m_um_v)$. There are no other edges.
Denote this edge-weighted, vertex-unweighted
graph by $G_1$. It has $N$ vertices, total edge weight at most one,
and weighted degree at most $2/N$ at each vertex.

A clone of $u$ has weighted degree
$\deg_{G^-}(u)/m_u\le2a_u/(Na_u)=2/N$.
The YES independent set becomes one of the same normalized size by
selecting all its clone blocks.

\begin{lemma}[Density after vertex unweighting]\label{lem:aks-clone-density}
In the NO case, every $S'\subseteq V(G_1)$ satisfies
\begin{equation}\label{eq:aks-clone-density}
 e_{G_1}(S')\ge\Gamma_\xi(|S'|/N)-\varepsilon_{\mathrm{dens}}-\varepsilon_{\mathrm{loop}}.
\end{equation}
\end{lemma}

\begin{proof}
Take any $S'\subseteq V(G_1)$ and let $f_u$ be the fraction of
clones of $u$ in $S'$. Include each original vertex $u$ in a random
set $S$ independently with probability $f_u$.
Since $G^-$ has no loops,
$\mathbb E a(S)=|S'|/N$ and
$\mathbb E e_{G^-}(S)=e_{G_1}(S')$.
Apply~\eqref{eq:aks-loop-free-density} and Jensen's inequality for
the convex function $\Gamma_\xi$.
This is the argument of~\cite[Section~4, Step~1]{AKS2011}, applied
after the loop removal of Lemma~\ref{lem:aks-loop-removal}.
\end{proof}

\subsubsection{Deterministic sparsification of a dense-set promise}

The next lemma converts this weighted density promise into the
required bounded-degree graph. We introduce the matrix notation
used in its proof. For a loopless edge-weighted graph $G$, its
\emph{Laplacian} $L_G$ is the symmetric matrix determined by
\[
 x^{\mathsf T}L_Gx=\sum_{uv\in E(G)}w_{uv}(x_u-x_v)^2
 \qquad(x\in\mathbb R^{V(G)}).
\]
The notation $L_G\preceq L_H$ means that
$x^{\mathsf T}L_Gx\le x^{\mathsf T}L_Hx$ for every vector $x$.
For the indicator vector $\mathbf1_S$, this quadratic form is the
weight of the cut of $S$. For $S=\{v\}$, it is the weighted degree
of $v$. A \emph{spectral sparsifier} uses fewer edges, with
possibly different positive weights, while preserving these
quadratic forms within specified multiplicative bounds.
Bounds on cuts and weighted degrees then control internal edge
weights through $2e_G(S)=\vol_G(S)-w(\cut_G(S))$.

\begin{lemma}\label{lem:deterministic-density-sparsification}
Fix positive constants $C_0,\zeta,\eta$. Let $G$ be a loopless
$N$-vertex edge-weighted graph of total edge weight at most one
and weighted degree at most $C_0/N$ at every vertex. Suppose that
every set of at least $\eta N$ vertices has internal edge weight
at least $\zeta$. There is a deterministic polynomial-time
construction of a simple unweighted graph $J$ on $V(G)$, using only
edges of $G$, whose
maximum degree is bounded by a constant depending only on
$C_0,\zeta$, and whose independent sets all have size less than
$\eta N$. Every independent set of $G$ remains independent in $J$.
\end{lemma}

\begin{proof}
Choose a fixed rational $0<\theta\le\min\{1/12,\zeta/14\}$,
and put $d=\lceil\theta^{-2}\rceil$.
The spectral sparsification theorem of Batson, Spielman, and
Srivastava~\cite[Theorem~1.1]{BSS2012}, followed by rescaling,
gives a graph $H$ on $V(G)$, using at most $dN$ of the edges of
$G$ with new positive weights, and satisfying
$L_G\preceq L_H\preceq(1+7\theta) L_G$.
Indeed, its distortion is at most
$((1+\theta)/(1-\theta))^2\le1+7\theta$.
Applying this inequality to $\mathbf1_{\{v\}}$ gives
$\deg_G(v)\le\deg_H(v)\le(1+7\theta) C_0/N$ and,
after summing, $\operatorname{vol}_H(S)\ge\operatorname{vol}_G(S)$.
Applying it to $\mathbf1_S$ gives
$w(\cut_H(S))\le(1+7\theta) w(\cut_G(S))$.
If $|S|\ge\eta N$, use
$2e_H(S)=\vol_H(S)-w(\cut_H(S))$ to obtain
\[
 e_H(S)\ge e_G(S)-\tfrac{7\theta}{2}w(\cut_G(S))
 \ge\zeta-\tfrac{7\theta}{2}\ge3\zeta/4,
\]
because every cut has weight at most one.

Discard every edge of $H$ whose weight is less than
$\zeta/(4dN)$, and let $J$ be the simple unweighted graph formed
by the remaining edges. Since $H$ has at most $dN$ edges, the
total weight discarded is at most $\zeta/4$. Every set of at
least $\eta N$ vertices therefore retains internal weight at
least $\zeta/2$, so it contains an edge of $J$. Thus every
independent set of $J$ has size less than $\eta N$.

Every retained edge incident with $v$ has weight at least
$\zeta/(4dN)$. Consequently,
$\deg_J(v)\le(4dN/\zeta)\deg_H(v)\le4(1+7\theta) C_0d/\zeta$.
This is a constant depending only on $C_0$ and $\zeta$.
Since $J$ uses only edges of $G$, every independent set of $G$
remains independent in $J$. The construction is deterministic
and polynomial time by the cited sparsification theorem.
\end{proof}

Combining the three steps gives the source promise used in the
UGC applications.

\begin{theorem}[Bounded-degree independent-set source]
\label{thm:bounded-degree-kr-source}
For every fixed rational $0<\eta<1/10$, there is a constant
$D_\eta$ such that the ordinary UGC gives a deterministic
polynomial-time hardness reduction to simple unweighted graphs
$G$ of maximum degree at most $D_\eta$, with
$\alphaI(G)\ge(1/2-\eta)|V(G)|$ in the YES case and
$\alphaI(G)<\eta|V(G)|$ in the NO case.
\end{theorem}

\begin{proof}
Set $p=1/2-\eta/4$ and $\xi=-p/(1-p)$. Fix a positive rational
$g$ with $4g<\Gamma_\xi(\eta)$. Use analytic error
$\varepsilon_{\mathrm{dens}}\le g$ and loop-loss budget
$\varepsilon_{\mathrm{loop}}=g$, and choose the Unique Games error sufficiently
small for AKS Theorem~3.1 and for $\gamma\le\eta/4$.
Pad the alphabet with a fixed $t$ satisfying
$(p^2+(1-p)^2)^t<g$. After dropping loops and replacing vertices by
copies, the YES independent-set density is at least
$p-\gamma\ge1/2-\eta/2$. In the NO case,
equation~\eqref{eq:aks-clone-density} and monotonicity show that
every set of at least $\eta N$ vertices has internal weight
greater than $\Gamma_\xi(\eta)-2g>2g$.
Apply Lemma~\ref{lem:deterministic-density-sparsification} with
$C_0=2$ and $\zeta=g$. Its output retains the YES
independent set and has no NO independent set of size $\eta N$.
Its degree bound depends only on $\eta$.
All alphabet parameters are constants for the fixed promise, so
the whole reduction is deterministic and polynomial time.
\end{proof}

\subsection{General sparsification theorems}
\label{subsec:general-sparsification}

We choose thresholds proportional to factor sizes. Bounding the
number of factors containing each vertex bounds both the total
completion loss and the degree after replacement, provided the
retained edges also have bounded degree.

Use the edge decomposition in~\eqref{eq:factor-edge-union} and the
replacement graph $J$ from Section~\ref{sec:bounded-degree-hardness}.
The input supplies $E_0$, the factor index sets, and each factor's
support or sides. Put $N=|V(H)|$.
Fix constants $d_0\ge0$,
$d_1\ge1$, and assume:
\begin{enumerate}
\item\label{hyp:retained-degree} $\Delta((V(H),E_0))\le d_0$.
\item\label{hyp:factor-incidence}
$|\{f\in\mathcal C\cup\mathcal B:v\in U_f\}|\le d_1$
for every $v\in V(H)$.
\end{enumerate}
Counting incidences between vertices and factor supports gives
\begin{equation}\label{eq:general-factor-incidence}
 \sum_{f\in\mathcal C\cup\mathcal B}|U_f|
 \overset{\text{Hyp.~\ref{hyp:factor-incidence}}}{\leq}d_1N.
\end{equation}

For a fixed real $\zeta>0$, the integer
$t_f=\lceil\zeta|U_f|\rceil\ge1$ is the number of vertices in
each set that the replacement of factor $f$ must join. Write $R_f$
for that replacement graph, and let
$\mathcal R=(R_f)_{f\in\mathcal C\cup\mathcal B}$ be the collection
of all replacements. Each $R_f$ is built on its prescribed support;
the output keeps $E_0$ and takes their union:
\[
 H\xrightarrow{\text{Lemma~\ref{lem:factor-replacement}}}(E_0,\mathcal R)
 \xrightarrow{\text{Lemma~\ref{lem:factor-replacement}}}J,\qquad
 V(J)=V(H),\qquad
 E(J)=E_0\cup\bigcup_{f\in\mathcal C\cup\mathcal B}E(R_f).
\]
The following lemma bounds the degree of this union and the sum of
the threshold losses.

\begin{lemma}[Common sparse factor replacements]
\label{lem:factor-replacement}
Let $H$ have the edge description~\eqref{eq:factor-edge-union},
and let fixed constants $d_0\ge0$, $d_1\ge1$ satisfy
conditions~\ref{hyp:retained-degree}--\ref{hyp:factor-incidence}.
Fix $\zeta>0$ and set
$t_f=\lceil\zeta|U_f|\rceil$ for every factor. There is a
deterministic polynomial-time construction of a spanning subgraph
$J\subseteq H$ with $E_0\subseteq E(J)$ that replaces each factor by a graph
satisfying its no-hole condition at threshold $t_f$, and has maximum degree at most
$d_0+d_1c_\zeta$, where $c_\zeta$ depends only on $\zeta$.
Moreover,
$\sum_{f\in\mathcal C\cup\mathcal B}(t_f-1)<\zeta d_1N$
whenever at least one factor is present.
\end{lemma}

\begin{proof}
Set $\kappa=\max\{1,1/\zeta-1\}$ and
$c_\zeta=d_{\zeta,\kappa}$ from Lemma~\ref{lem:factor-expanders}.
Use that lemma for each clique factor. For a biclique with side
sizes $a,b$, put $u=\lceil\zeta(a+b)\rceil$. If $\min\{a,b\}<u$,
use the empty graph on its support: no admissible pair of $u$-sets
exists. Otherwise $\min\{a,b\}\ge u\ge\zeta(a+b)$, so
\[
 \frac{\max\{a,b\}}{\min\{a,b\}}
 =\frac{a+b}{\min\{a,b\}}-1
 \le\frac1\zeta-1\le\kappa.
\]
Lemma~\ref{lem:factor-expanders} therefore applies to this factor.
Each replacement has degree at most $c_\zeta$ and satisfies the
no-hole condition at threshold $t_f$. If $\zeta>1/2$ the nonempty
branch is impossible, and if $\zeta=1/2$
it requires $a=b$.

Take the union of the replacements with $E_0$. The replacements of
vacuous biclique factors need not be connected. The completion
lemmas require only the no-hole conditions; when connected factor
replacements are needed separately, use Lemma~\ref{lem:factor-expanders}
with its side-ratio assumption.
For every $v\in V(H)$, the degree in the union satisfies
\[
 \deg_J(v)\overset{\text{Lem.~\ref{lem:factor-expanders}}}{\leq}
 \deg_{(V(H),E_0)}(v)
 +|\{f\in\mathcal C\cup\mathcal B:v\in U_f\}|c_\zeta
 \overset{\text{Hyps.~\ref{hyp:retained-degree}, \ref{hyp:factor-incidence}}}{\leq}
 d_0+d_1c_\zeta.
\]
Finally, the ceiling defining $t_f$ gives
\begin{equation}\label{eq:factor-threshold-slack}
 t_f-1<\zeta|U_f|\qquad(f\in\mathcal C\cup\mathcal B).
\end{equation}
If at least one factor is present, then
\[
 \sum_{f\in\mathcal C\cup\mathcal B}(t_f-1)
 \overset{\eqref{eq:factor-threshold-slack}}{<}
 \zeta\sum_{f\in\mathcal C\cup\mathcal B}|U_f|
 \overset{\eqref{eq:general-factor-incidence}}{\leq}\zeta d_1N.
\]
Each factor has nonempty support, so there are at most
$d_1N$ factors. Let $h\ge1$ be the fixed walk length chosen in
Appendix~\ref{app:factor-expanders} for this $\zeta$. The construction
enumerates $O(8^h\sum_{f\in\mathcal C\cup\mathcal B}|U_f|)
=O(8^h d_1N)$ auxiliary steps.
Sorting the resulting endpoint pairs and merging repetitions uses
$O(8^h d_1N\log(2+8^h d_1N))$ comparisons, so the construction
is deterministic and polynomial time for the fixed parameters.
\end{proof}

\begin{theorem}[Treedepth sparsification]
\label{thm:td-sparsification}
Let $H$ have the edge description~\eqref{eq:factor-edge-union},
and let fixed constants $d_0\ge0$, $d_1\ge1$ satisfy
conditions~\ref{hyp:retained-degree}--\ref{hyp:factor-incidence}.
For every fixed $\varepsilon>0$
one can construct deterministically in polynomial time a spanning
subgraph $J\subseteq H$ with $E_0\subseteq E(J)$ whose maximum degree depends only on
$d_0,d_1,\varepsilon$ and such that
$\td(H)-\varepsilon N\le\td(J)\le\td(H)$.
Every elimination forest of $J$ can be transformed into one of
$H$ whose height increases by at most $\varepsilon N$.
\end{theorem}

\begin{proof}
Apply Lemma~\ref{lem:factor-replacement} with
$\zeta=\varepsilon/(4d_1)$. Its degree bound depends only on
the stated constants. The exception bound in
Lemma~\ref{lem:factor-sparsification} satisfies
\begin{equation}\label{eq:common-forest-exception-budget}
 \operatorname{exc}_{\mathrm{td}}
 \overset{\eqref{eq:completion-deletion-bounds}}{=}
 2\sum_{f\in\mathcal C}(t_f-1)+4\sum_{f\in\mathcal B}(t_f-1) \le4\sum_{f\in\mathcal C\cup\mathcal B}(t_f-1)
 \overset{\text{Lem.~\ref{lem:factor-replacement}}}{\leq}
 4\zeta d_1N=\varepsilon N.
\end{equation}
Lemma~\ref{lem:factor-sparsification} now gives both conclusions.
\end{proof}

\begin{theorem}[Treewidth sparsification]
\label{thm:tw-sparsification}
Let $H$ have the edge description~\eqref{eq:factor-edge-union},
and let fixed constants $d_0\ge0$, $d_1\ge1$ satisfy
conditions~\ref{hyp:retained-degree}--\ref{hyp:factor-incidence}.
Fix an integer $s\ge1$. Suppose the supplied factor description also
includes, for every $f\in\mathcal B$, integers
$1\le\ell_f,m_f\le s$ and partitions
$A_f=A_{f,1}\mathbin{\dot\cup}\cdots\mathbin{\dot\cup}A_{f,\ell_f}$
and $B_f=B_{f,1}\mathbin{\dot\cup}\cdots\mathbin{\dot\cup}B_{f,m_f}$
into nonempty sets, together with indices
$g_{f,i},h_{f,j}\in\mathcal C$ such that
$A_{f,i}\subseteq U_{g_{f,i}}$ and $B_{f,j}\subseteq U_{h_{f,j}}$
for every $i\in[\ell_f]$ and $j\in[m_f]$.
For every fixed $\varepsilon>0$, one can construct deterministically
in polynomial time a spanning subgraph $J\subseteq H$ with
$E_0\subseteq E(J)$ whose maximum
degree depends only on $d_0,d_1,s,\varepsilon$ and such that
$\operatorname{tw}(H)-\varepsilon N
 \le\operatorname{tw}(J)\le\operatorname{tw}(H)$.
Every tree decomposition of $J$ can be transformed into one of
$H$ whose width increases by at most $\varepsilon N$.
\end{theorem}

\begin{proof}
Apply Lemma~\ref{lem:factor-replacement} to the original factors
with $\zeta=\varepsilon/(d_1\max\{3,2s^2\})$. It constructs $J$
with maximum degree depending only on the stated constants and
thresholds $t_f=\lceil\zeta|U_f|\rceil$.

For $f\in\mathcal B$, the supplied partitions and their indices give
\begin{equation}\label{eq:tw-factor-piece-count}
 \ell_fm_f\leq s^2,\qquad
 A_f\times B_f=\mathop{\dot\bigcup}_{i\in[\ell_f],\,j\in[m_f]}
 (A_{f,i}\times B_{f,j}).
\end{equation}
Describe the same graph $H$ by keeping $E_0$ and its clique
factors and replacing the biclique indexed by $f$ with the
$\ell_fm_f$ bicliques indexed by $(f,i,j)$, on sides
$A_{f,i},B_{f,j}$. Give each of these bicliques threshold $t_f$.
Each has its sides contained in the clique supports $U_{g_{f,i}}$
and $U_{h_{f,j}}$.

The same graph $J$ satisfies this refined factor description. The
clique conditions hold by Lemma~\ref{lem:factor-replacement}.
Any $t_f$-subsets of $A_{f,i}$ and $B_{f,j}$ are also subsets of
$A_f$ and $B_f$, so the original replacement joins them. If either
new side has fewer than $t_f$ vertices, this condition is vacuous.

Lemma~\ref{lem:tw-factor-sparsification} charges at most $3(t_f-1)$
per clique and $2(t_f-1)$ per refined biclique. Thus its completion
increases width by at most
\[
\begin{aligned}
 &3\sum_{f\in\mathcal C}(t_f-1)
 +2\sum_{f\in\mathcal B}\ell_fm_f(t_f-1)
 \quad\overset{\eqref{eq:tw-factor-piece-count}}{\leq}
 3\sum_{f\in\mathcal C}(t_f-1)
       +2s^2\sum_{f\in\mathcal B}(t_f-1)\\
 &\quad\le\max\{3,2s^2\}\sum_{f\in\mathcal C\cup\mathcal B}(t_f-1)
 \quad\overset{\text{Lem.~\ref{lem:factor-replacement}}}{\leq}
 \max\{3,2s^2\}\zeta d_1N=\varepsilon N.
\end{aligned}
\]
The comparison lemma also gives the stated parameter inequalities.
\end{proof}

\Needspace{9\baselineskip}
\begin{theorem}[Total-depth sparsification]
\label{thm:total-depth-sparsification}
Let $H$ have the edge description~\eqref{eq:factor-edge-union},
and let fixed constants $d_0\ge0$, $d_1\ge1$ satisfy
conditions~\ref{hyp:retained-degree}--\ref{hyp:factor-incidence}.
For every fixed $\varepsilon>0$
one can construct deterministically in polynomial time a spanning
subgraph $J\subseteq H$ with $E_0\subseteq E(J)$ whose maximum degree depends only on
$d_0,d_1,\varepsilon$ and such that
$\totd(H)-\varepsilon N^2\le\totd(J)\le\totd(H)$.
Every elimination forest $F$ of $J$ can be transformed into an
elimination forest $\widehat F$ of $H$ satisfying
$\totd(\widehat F)\le\totd(F)+\varepsilon N^2$.
\end{theorem}

\begin{proof}
Use the same factor replacements as in
Theorem~\ref{thm:td-sparsification}, with
$\zeta=\varepsilon/(4d_1)$. Equation~\eqref{eq:common-forest-exception-budget}
gives $\operatorname{exc}_{\mathrm{td}}\le\varepsilon N$, and
Lemma~\ref{lem:total-depth-comparison} gives the completion and both
parameter bounds, with total-depth increase at most
$N\operatorname{exc}_{\mathrm{td}}
\overset{\eqref{eq:common-forest-exception-budget}}{\leq}\varepsilon N^2$.
\end{proof}

If a construction has a fixed number $d$ of designated vertices
that are not needed for sparsification, deleting them lowers treedepth and treewidth by
at most $d$. Restoring them above all existing roots increases total
depth by at most $dN$, where $N$ is the original graph order.
These costs are $O(1)$ for treedepth and treewidth, whose promise
bounds are linear in $N$, and $O(N)$ for total depth, whose bounds
are quadratic. They are therefore $o(Y)$ for the corresponding YES
upper bound $Y$, so
Lemma~\ref{lem:gap-preservation} absorbs them on sufficiently large
instances. We use this to omit the three vertices
$z_A,z_B,z_C$ from the NP treedepth construction.


\Needspace{10\baselineskip}
\subsection{Hardness consequences}
\label{subsec:bounded-degree-consequences}

Fix a target factor $r$ below a dense hardness threshold, and choose
a dense factor $r'>r$. Write $N$ for the dense graph's order and
$Y$ for its YES upper bound. It is linear
in $N$ for treedepth and treewidth and quadratic for total depth.
For $\tau>0$, the sparsification losses $\tau N$ and $\tau N^2$ can therefore be made
at most $\delta Y$ by choosing $\tau$ small. Taking
$0<\delta<(r'-r)/(r+1)$, Lemma~\ref{lem:gap-preservation} transfers
the factor-$r$ hardness.
We verify the factor conditions below, using the following observation
for the SSE construction.
In these applications we choose connected replacements using
Lemma~\ref{lem:factor-expanders}. The biclique side-size ratios,
when bicliques occur, are bounded as checked below. Although the
general sparsification theorems do not require balance, these
connected replacements preserve the connected components of the
dense graph: each replaced factor still connects its entire support.

The transfer steps are summarized below. Here $H$ is a dense graph,
$H'$ its sparse replacement, $R$ its regular extension, and
$R_\ell$ replaces each edge of $R$ by a path of fixed even length
$\ell$. Sparsification and regularization apply to all three objectives;
subdivision applies to treedepth and treewidth, and to total depth under SSE.
The guarantees are Theorems~\ref{thm:td-sparsification}--\ref{thm:total-depth-sparsification},
Proposition~\ref{prop:regularization}, and
Corollary~\ref{cor:subdivision-hardness}, respectively:
\[
 H\xrightarrow{\text{sparsify}}H'
 \xrightarrow{\text{regularize}}R
 \xrightarrow{\text{subdivide}}R_\ell.
\]

\begin{lemma}[Clique factors of a line graph]\label{lem:linegraph-factors}
Let $G$ be an unweighted host multigraph. Identify each edge
occurrence $e\in E(G)$ with its vertex in $\mathcal L(G)$, and put
$U_v=\{e\in E(G):v\text{ is an endpoint of }e\}$ for $v\in V(G)$.
The nonempty sets
$U_v$ give clique factors for $\mathcal L(G)$, and each line-graph
vertex belongs to at most two factors.
\end{lemma}

\begin{proof}
Two line-graph vertices are adjacent if and only if their host edges
share an endpoint, so
$E(\mathcal L(G))=\bigcup_{v\in V(G)}\binom{U_v}{2}$.
An edge has at most two endpoints. A loop contributes one vertex
to the factor at its endpoint.
\end{proof}

We may require arbitrarily large sources: disjoint copies preserve the satisfiable fraction
of a parity system and the independence fraction and maximum degree
of an independent-set source. A relative independent-set
approximation on the copies transfers to the original by choosing
the copy with the largest selected set. For SSE, parallel copies
of every host edge preserve normalized volumes and normalized cut sizes and
multiply the line-graph order.

In all SSE cases, let $c>0$ and $\varepsilon_0>0$ be the constants
in Lemma~\ref{lem:sse-unweighted-gap}. Apply it with sufficiently
small source error $0<\eta<\varepsilon_0$ and
$q=\lceil2/\eta\rceil$, as in Section~\ref{subsec:sse-dense-gap}.
Let $G^*$ denote its host, put $H=\mathcal L(G^*)$, and write
$N=|E(G^*)|=|V(H)|$, counting edge multiplicity, and
$\beta=(c/2)\sqrt\eta\le1/4$.
The line graph has YES height bound $4\eta N$ and NO balanced-separator
bound $\beta N$, the latter by
Lemma~\ref{lem:linegraph-separators} with $\rho=1/2$.
Lemma~\ref{lem:linegraph-factors} gives $d_1=2$, $d_0=0$, and no
biclique factors, regardless of the host degree. Sparsification
replaces these host-incidence cliques directly; the sparse output
need not itself be a line graph. Since $\beta/\eta$ grows
without bound as $\eta$ decreases, any fixed target gap can be
obtained before sparsification.

\Needspace{7\baselineskip}
\begin{corollary}[Bounded-degree treedepth hardness]
\label{cor:td-bounded-degree}
\label{thm:np-bounded-degree}
\label{thm:ugc-bounded-degree}
\label{thm:sse-bounded-degree}
For each fixed $r\ge1$, approximating treedepth within factor $r$
on simple unweighted graphs of bounded maximum degree is
NP-hard if $r<18/17$, Unique Games-hard if $r<4/3$, and
SSE-hard for every fixed $r$.
The degree bound depends only on the excluded factor. All reductions
are deterministic.
\end{corollary}

\begin{proof}
Apply Theorem~\ref{thm:td-sparsification} and
Lemma~\ref{lem:gap-preservation} in each case.

\emph{NP-hardness.}
Retain every edge of $G_F$, whose maximum degree is at most
$3b_\epsilon$ by Lemma~\ref{lem:xor-cover}. After omitting
$z_A,z_B,z_C$, replace the three anchor cliques and the three
anchor--part bicliques. Each vertex lies in at most two factor
supports, and each biclique has both sides of size $3m$.
Thus $d_1=2$, $d_0=3b_\epsilon$, and $\kappa=1$.
Deleting the three designated vertices costs at most $3$ in treedepth.
The remaining order is $18m$, so sparsification costs at most $18\tau m$.
The dense gap has positive linear slack; choose $\tau$ small and then
$m$ large. Appendix~\ref{app:transfer-td-np} gives the exact bounds
and choices.

\emph{UGC hardness.}
Fix a target factor $r<4/3$, and choose $r'$ with
$r<r'<4/3$. Take $q,\eta$ from the dense gap calculation in
Section~\ref{subsec:ugc-td-gap} for factor $r'$, and apply
Theorem~\ref{thm:bounded-degree-kr-source} with this $\eta$ in
$S_q(G)$, where $n=|V(G)|$. Replace the $q$ layer cliques by
bounded-degree expanders and retain all cross-layer source edges. These retained edges have
maximum degree at most $(q-1)d_\eta$, so $d_1=1$ and
$d_0=(q-1)d_\eta$. A layer clique gives a YES scale of at least $n$, while the graph
has $qn$ vertices. Thus the loss is at most $q\tau Y$ and can be
absorbed into the factor slack; see Appendix~\ref{app:transfer-td-ugc}.

\emph{SSE hardness.}
Use the common line-graph factors above.
Lemma~\ref{lem:balanced-by-td} converts the separator bound into
NO treedepth bound $\beta N$. Since $\beta/\eta$ is unbounded as $\eta$ tends to zero, first
choose a dense gap larger than the target factor, then make $\tau$
small enough to preserve it. The calculation is in
Appendix~\ref{app:transfer-td-sse}.

\end{proof}

The NP-hardness assertion answers the bounded-degree conjecture of
Schirrmacher, Siebertz, and
Vigny~\cite[Conjecture~1]{SchirrmacherSiebertzVigny2025}.
Weighted treedepth hardness at maximum degree ten is proved
in~\cite{DirksEtAl2025Weighted}.

\Needspace{7\baselineskip}
\begin{corollary}[Bounded-degree treewidth hardness]
\label{cor:tw-bounded-degree}
\label{thm:tw-bounded-degree}
\label{thm:tw-ugc-bounded-degree}
\label{thm:tw-sse-bounded-degree}
For each fixed $r\ge1$, approximating treewidth within factor $r$
on simple unweighted graphs of bounded maximum degree is
NP-hard if $r<12/11$, Unique Games-hard if $r<4/3$, and
SSE-hard for every fixed $r$.
The degree bound depends only on the excluded factor. All reductions
are deterministic.
\end{corollary}

\begin{proof}
Apply Theorem~\ref{thm:tw-sparsification} and
Lemma~\ref{lem:gap-preservation}.

\emph{NP-hardness.}
In $K_R(G_F)$, replace the three global cliques and retain every
edge of each source copy. Lemma~\ref{lem:xor-cover} bounds the
retained degree by $d_0=3b_\epsilon$. Each vertex lies in one
clique factor, so $d_1=1$. There are no biclique factors.
The clique replacements use $\kappa=1$ in
Lemma~\ref{lem:factor-expanders}.
The dense gap has positive slack proportional to $m$, and the
sparsification loss is $9\tau Rm$. A sufficiently small fixed $\tau$
preserves the gap; Appendix~\ref{app:transfer-tw-np} records its
coefficient and the exact choice.

\emph{UGC hardness.}
Fix a target factor $r<4/3$, and choose $r'$ with
$r<r'<4/3$. Take $\zeta,R$ from the dense gap calculation in
Section~\ref{subsec:ugc-tw-gap} for factor $r'$, and use
Theorem~\ref{thm:bounded-degree-kr-source} with $\eta=\zeta/2$,
as in Theorem~\ref{thm:ugc-tripartite-cover}. This is admissible
since $0<\eta<1/24<1/10$. Apply it in
$K_R(P_3(G))$, where $n=|V(G)|$.
Replace only the three global cliques by bounded-degree expanders
and retain every product edge in each source copy. These source
edges already have bounded degree, so
$d_1=1$ and $d_0=2d_\eta=D_{3,\zeta}$, with no biclique factors.
Lemma~\ref{lem:multipartite-product} gives YES cover bound
$(3/2+3\zeta/2)n\leq(3/2+3\zeta)n$ and NO cover bound
$(2-\zeta)n\geq(2-3\zeta)n$ for $P_3(G)$, as required by
Section~\ref{subsec:ugc-tw-gap}.
The global cliques give YES scale at least $Rn-1$, compared with
order $3Rn$. A small fixed $\tau$ therefore absorbs the loss into
the factor slack; see Appendix~\ref{app:transfer-tw-ugc}.

\emph{SSE hardness.}
Use the common line-graph factors.
Lemma~\ref{lem:treewidth-balanced-separator} gives
$b(H)\leq\operatorname{tw}(H)+1$ for the dense line graph.
Thus the sparse NO width is at least $\beta N-1-\tau N$.
The YES line-graph construction has treedepth at most $4\eta N$.
An elimination forest of height $h$ gives a tree decomposition of
width at most $h-1$: connect its roots and use its ancestor paths as
bags. A fixed vertex occurs on its descendant subtree, and every
edge is covered at its deeper endpoint. Hence the dense YES
treewidth is at most $4\eta N$. By
Lemma~\ref{lem:minor-monotonicity}, deleting edges to obtain the sparse
graph cannot increase this bound.
Choose $\eta$ for positive linear gap slack, then $\tau$ small and
$N$ large enough to absorb the sparsification loss and the constant
$1$. Appendix~\ref{app:transfer-tw-sse} gives the choices.

\end{proof}

The SSE treewidth hardness at every fixed factor on unrestricted
graphs is due to Wu, Austrin, Pitassi, and Liu~\cite[Theorem~1.1]{APW}. The bounded-degree
statement here follows by applying our treewidth sparsification
theorem to the line-graph construction.

For the total-depth transfer below, we use the following bound for
lifting an approximate solution through a reduction.

\begin{lemma}[Lifting approximate solutions]
\label{lem:solution-lifting}
Suppose each solution of a minimization instance $I'$ with cost $z$
can be converted in polynomial time into a solution of $I$ with cost
at most $z/a+\eta L(I)$, where $a>0$, $\eta\ge0$, and the bound $L(I)$ satisfies
\begin{equation}\label{eq:solution-lifting-lower}
 0\le L(I)\le\OPT(I).
\end{equation}
If
\begin{equation}\label{eq:solution-lifting-upper}
 \OPT(I')\le a\OPT(I)+\eta aL(I),
\end{equation}
then, for every $r\ge1$, an $r$-approximation for
$I'$ gives an $(r+(r+1)\eta)$-approximation for $I$. If instead
\begin{equation}\label{eq:solution-lifting-exact-upper}
 \OPT(I')\le a\OPT(I),
\end{equation}
the factor is $r+\eta$.
\end{lemma}

\begin{proof}
An $r$-approximate solution has $z\le r\OPT(I')$, so its conversion
has cost at most
\[
\frac za+\eta L(I)
 \le\frac{r\OPT(I')}{a}+\eta L(I) \overset{\eqref{eq:solution-lifting-upper}}{\le}
 r\OPT(I)+(r+1)\eta L(I)
 \overset{\eqref{eq:solution-lifting-lower}}{\le}
 \bigl(r+(r+1)\eta\bigr)\OPT(I).
\]
Under~\eqref{eq:solution-lifting-exact-upper}, the conversion cost satisfies
\[
 \frac za+\eta L(I)
 \overset{\eqref{eq:solution-lifting-exact-upper}}{\le}r\OPT(I)+\eta L(I)
 \overset{\eqref{eq:solution-lifting-lower}}{\le}(r+\eta)\OPT(I).
\]
\end{proof}

\Needspace{8\baselineskip}
\begin{corollary}[Bounded-degree total-depth hardness]
\label{cor:total-depth-bounded-degree}
\label{thm:np-total-depth-bounded-degree}
\label{thm:total-depth-ugc-bounded-degree}
\label{cor:total-depth-sse-bounded-degree}
Approximating total depth within every fixed factor $1\le r<63/62$
is NP-hard on simple unweighted graphs of bounded maximum degree.
On such graphs,
total depth is also Unique Games-hard to approximate within every
fixed factor $1\le r<8/7$, and SSE-hard to approximate within every
fixed factor $r\ge1$. The degree bound depends only on the excluded
factor. All reductions are deterministic.
\end{corollary}

\begin{proof}
Apply Theorem~\ref{thm:total-depth-sparsification}.
It loses at most $\tau N^2$ and completes sparse forests at the
same additive cost.

\emph{NP-hardness.}
Fix $1\le r<r'<63/62$. Use the dense family
$H=H^\Sigma_M(G_F^\Sigma)$ from
Theorem~\ref{thm:np-total-depth-dense-gap} for factor $r'$, where
$F$ is a parity system with $m$ equations and occurrence bound
$b_\epsilon$ from Theorem~\ref{thm:xor-source}.
The three part cliques and the three bicliques $X\times T_X$
for $X\in\{A,B,C\}$
give $d_1=2$. Retain every edge of $G_F^\Sigma$; its maximum degree
is at most $14b_\epsilon$ by~\eqref{eq:np-total-depth-degree}, so
$d_0=14b_\epsilon$. The part sizes are $27m,36m,36m$ by
\eqref{eq:np-total-depth-part-sizes}, while each $T_X$ has size
$M=\Theta(m)$. Thus all biclique side-size ratios are bounded by
a constant depending only on $r'$.

The dense family has $N=99m+3M=\Theta(m)$ and total depth at
least $Mm=\Omega(N^2)$. Choose $\tau$ small relative to this
quadratic scale and the gap slack. This preserves the gap and also
ensures $\totd(H')\ge Mm/2$ for the sparse graph $H'$, as proved
in Appendix~\ref{app:transfer-total-np}. The latter bound will be
used for regularization.

\emph{UGC hardness.}
Fix a target factor $r<8/7$, and choose $r'$ with
$r<r'<8/7$. Use the parameters $q,\eta$ from the dense gap
calculation in Section~\ref{subsec:ugc-total-depth-gap} for factor $r'$, starting from
Theorem~\ref{thm:bounded-degree-kr-source}. Replace the $q$ layer
cliques by bounded-degree expanders and retain every cross-layer
source edge. Thus $d_1=1$ and $d_0=(q-1)d_\eta$.
Write $n=|V(G)|$ and $N=qn$. The layer cliques force total depth at least
$q n(n+1)/2\ge N^2/(2q)$. Hence the loss $\tau N^2$ is a
small fraction of the YES scale when $\tau$ is small. See
Appendix~\ref{app:transfer-total-ugc} for the factor calculation.

\emph{SSE hardness.}
The YES line-graph construction has height at most $4\eta N$ and
has $N$ vertices, so its total depth is at most
$N\cdot4\eta N=4\eta N^2$. Lemma~\ref{lem:balanced-by-td} gives NO total depth
at least $\beta N^2/2$. The ratio of these quadratic bounds is unbounded as $\eta$ tends
to zero. Choose the dense gap first and then $\tau$ small enough
to absorb the quadratic loss; Appendix~\ref{app:transfer-total-sse}
gives the explicit slack.

\end{proof}

\subsection{Regular graphs and subdivisions}

Both constructions below recover a forest by contracting connected
vertex sets. The following observation records the common argument.

\begin{lemma}[Projection through connected parts]
\label{lem:connected-part-projection}
Let $G,H$ be finite simple graphs, and let
$(X_v)_{v\in V(G)}$ partition $V(H)$ into nonempty connected sets.
Suppose every edge $uv\in E(G)$ has an edge of $H$ between $X_u$
and $X_v$. From any elimination forest $F$ of $H$, one can compute
in polynomial time an elimination forest $\overline F$ of $G$ such that
\[
 \operatorname{depth}_{\overline F}(v)\le\operatorname{depth}_F(x)
 \quad(v\in V(G),\ x\in X_v).
\]
Consequently $h(\overline F)\le h(F)$ and
$\totd(F)\ge\sum_{v\in V(G)}|X_v|\operatorname{depth}_{\overline F}(v)$.
\end{lemma}

\begin{proof}
For each $v$, choose the shallowest vertex $r_v$ of $X_v$.
Lemma~\ref{lem:connected-set-ancestor} makes it unique and an ancestor
of every vertex in $X_v$. If $xy$ witnesses an edge $uv\in E(G)$,
with $x\in X_u$ and $y\in X_v$, then $x,y$ are comparable in $F$.
Both $r_u,r_v$ are ancestors of their deeper endpoint, so they too
are comparable. Restrict $F$ to the representatives and rename
$r_v$ as $v$; this gives $\overline F$. Restriction and the ancestor
property give, for every $x\in X_v$,
\[
 \operatorname{depth}_{\overline F}(v)
 \le\operatorname{depth}_F(r_v)
 \overset{\text{Lem.~\ref{lem:connected-set-ancestor}}}{\le}
 \operatorname{depth}_F(x).
\]
Taking the maximum and summing over the partition prove the two
consequences. Computing depths, choosing representatives, and restricting
the forest are polynomial-time operations. No connectedness of $G$ or
$H$ is required; for empty graphs the empty partition and forest suffice.
\end{proof}

\begin{proposition}\label{prop:regularization}
For every finite simple graph $G$ of order $n$ and fixed odd $d\geq3$ with
$d>\Delta(G)$, one can construct in polynomial time a simple
$d$-regular graph $R$ containing $G$ as an induced subgraph, with
$(d+2)n\leq |V(R)|\leq(d+3)n$ and
\[
\begin{gathered}
 \td(G)\leq\td(R)\leq\td(G)+d+2,\qquad \operatorname{tw}(G)\leq\operatorname{tw}(R)
 \leq\max\{\operatorname{tw}(G),d+2\},\\
 (d+2)\totd(G)\leq\totd(R)
 \leq(d+3)\totd(G)+\tfrac12(d+2)(d+3)n.
\end{gathered}
\]
\end{proposition}

\begin{proof}
For each $v\in V(G)$, set $t=d-\deg_G(v)$, so $1\le t\le d$.
Take graphs $Q_v$ whose vertex sets are pairwise disjoint and
disjoint from $V(G)$. Choose $P_v\subseteq V(Q_v)$ with $|P_v|=t$,
and join $v$ to every vertex of $P_v$. Define $Q_v$ as follows.
\begin{enumerate}
 \item If $t$ is even, take $K_{d+1}$ and remove a perfect matching
   on $P_v$.
 \item If $t\ge3$ is odd, take $K_{d+2}$ and remove a cycle with
   vertex set $P_v$ and a perfect matching on $V(Q_v)\setminus P_v$.
   The latter set has even size $d+2-t$.
 \item If $t=1$, write $P_v=\{p\}$, take $K_{d+2}$, and choose
   distinct $a,b\in V(Q_v)\setminus\{p\}$. Remove $pa,pb$ and a
   perfect matching on $V(Q_v)\setminus\{p,a,b\}$.
\end{enumerate}
Define the graph $R$ by
\[
V(R)=V(G)\mathbin{\dot\cup}
       \mathop{\dot\bigcup}_{v\in V(G)}V(Q_v),\qquad E(R)=E(G)\cup\bigcup_{v\in V(G)}E(Q_v)
       \cup\bigcup_{v\in V(G)}\{\{v,p\}:p\in P_v\}.
\]
The three cases give
\[
 \deg_{Q_v}(p)=d-1\quad(p\in P_v),\qquad
 \deg_{Q_v}(x)=d\quad(x\in V(Q_v)\setminus P_v).
\]
Hence $\deg_R(v)=\deg_G(v)+|P_v|=d$ for every $v\in V(G)$,
and $\deg_R(x)=d$ for every $x\in V(Q_v)$.
Each $Q_v$ has minimum degree $d-1$, and its order satisfies
\begin{equation}\label{eq:regularization-attachment-size}
 d+1\le |V(Q_v)|\le d+2\qquad(v\in V(G)).
\end{equation}
Every connected component therefore has at least $d$ vertices,
so $Q_v$ is connected because
$2d>d+2$ for $d\geq3$. The resulting graph $R$ is simple.

Given an elimination forest $F$ of $G$, appending $V(Q_v)$ as a chain
below $v$ increases the height by at
most $d+2$. If the chain has $|V(Q_v)|$ vertices, their total depth is
at most
\[
 |V(Q_v)|\operatorname{depth}_F(v)
 +\frac{|V(Q_v)|(|V(Q_v)|+1)}2
 \overset{\eqref{eq:regularization-attachment-size}}{\leq}
 (d+2)\operatorname{depth}_F(v)+\frac{(d+2)(d+3)}2.
\]
Summing over $v$ and including the original vertices gives the stated
total-depth upper bound. Attach a bag $\{v\}\cup V(Q_v)$ to an
existing bag containing $v$ in a tree decomposition of $G$.
Doing this for every $v$ proves the treewidth bound.

For the total-depth lower bound, partition $V(R)$ into connected sets
$X_v=\{v\}\cup V(Q_v)$. Each edge $uv$ of $G$ itself witnesses
an edge between $X_u$ and $X_v$. Thus any elimination forest $F$
of $R$ projects by Lemma~\ref{lem:connected-part-projection} to a
forest $\overline F$ of $G$ with
\begin{equation}\label{eq:regularization-depth-projection}
 \operatorname{depth}_F(x)
 \overset{\text{Lem.~\ref{lem:connected-part-projection}}}{\ge}
 \operatorname{depth}_{\overline F}(v)
 \qquad(v\in V(G),\ x\in X_v).
\end{equation}
Consequently,
\[
 \totd(F)
 \overset{\eqref{eq:regularization-depth-projection}}{\geq}
 \sum_{v\in V(G)}|X_v|\operatorname{depth}_{\overline F}(v)
 \overset{\eqref{eq:regularization-attachment-size}}{\geq}
 (d+2)\totd(\overline F)\geq(d+2)\totd(G).
\]
This argument also gives a polynomial-time projection of forests.
\end{proof}

\begin{corollary}\label{cor:regular-hardness}
The hardness ranges in Corollaries~\ref{cor:td-bounded-degree},
\ref{cor:tw-bounded-degree}, and~\ref{cor:total-depth-bounded-degree}
also hold on $d$-regular graphs for a suitable fixed odd degree $d$
depending only on the excluded factor.
\end{corollary}

\begin{proof}
Choose a source gap larger than the target gap and apply
Proposition~\ref{prop:regularization}. Its fixed additive costs for
treedepth and treewidth vanish relative to their linear YES scale.
For total depth, let $G$ be the source graph, put $n=|V(G)|$, and
divide costs by $d+2$: the error is
$O(n^2/d+dn)$. Choose $d$ large, then $n$ large, relative to the
quadratic scale. In the NP case, with $M,m$ from the NP construction,
the bound $L=Mm/2\le\totd(G)$
from Appendix~\ref{app:transfer-total-np} and the proposition's
forest projection give the same approximation transfer directly.
Appendix~\ref{app:transfer-regular} supplies all parameter choices
and the projection calculation.
\end{proof}

\begin{corollary}\label{cor:subdivision-hardness}
For every fixed $g\geq3$, the treedepth and treewidth hardness ranges above
hold on bounded-degree graphs that are bipartite, $2$-degenerate, and
have girth at least $g$. Here $2$-degenerate means that every
nonempty subgraph has a vertex of degree at most two, and girth is the
length of a shortest cycle (infinite for an acyclic graph). On the same class, total depth is
SSE-hard to approximate within every fixed factor.
\end{corollary}

\begin{proof}
Replace every edge of $G$ by a path of fixed even length
$\ell\geq\max\{2,g/3\}$, obtaining $G_\ell$. All original vertices
lie in one side of a bipartition, since every replacement path has
even length. Each internal vertex has degree two, and deleting all
internal vertices leaves the original vertices isolated, so
$G_\ell$ is $2$-degenerate. Original vertices retain their degrees,
so $\Delta(G_\ell)\leq\max\{\Delta(G),2\}$.
Every cycle uses at least three
replacement paths and therefore has length at least $3\ell\geq g$.
By Lemma~\ref{lem:minor-monotonicity}, contracting each replacement
path gives
$\operatorname{tw}(G)
\overset{\text{Lem.~\ref{lem:minor-monotonicity}}}{\leq}
\operatorname{tw}(G_\ell)$ and
$\td(G)\overset{\text{Lem.~\ref{lem:minor-monotonicity}}}{\leq}
\td(G_\ell)$. Conversely, for each path
$u=x_0,x_1,\ldots,x_\ell=v$, attach the bags
$\{u,v,x_1\},\{v,x_1,x_2\},\ldots,
\{v,x_{\ell-2},x_{\ell-1}\}$ along a path to an old bag containing
$u,v$. These bags cover the replacement path and preserve the
connectedness condition for each vertex, so
$\operatorname{tw}(G_\ell)\leq\max\{\operatorname{tw}(G),2\}$.
If $\operatorname{tw}(G)\ge2$, the two bounds give equality.
If $\operatorname{tw}(G)\le1$, then $G$ is a forest. Subdivision
preserves acyclicity and whether an edge exists, so it preserves
treewidth zero or one as well. Thus treewidth is preserved in every case.
A path on $s$ vertices has an
elimination tree of height at most $\lceil\log_2(s+1)\rceil$: choose
a middle vertex as root and recurse on the two remaining subpaths.
Thus the $\ell-1$ internal vertices on each replaced edge have such
a tree of height $h=\lceil\log_2\ell\rceil$. Attach this tree below
the deeper of the original edge's endpoints in an elimination forest
of $G$. The original endpoints are comparable, so the edge to the
shallower endpoint is represented as well. This gives
\begin{equation}\label{eq:subdivision-height-upper}
 \td(G_\ell)\leq\td(G)+h,\qquad h=\lceil\log_2\ell\rceil.
\end{equation}
Lemma~\ref{lem:gap-preservation} absorbs the fixed height increase
in~\eqref{eq:subdivision-height-upper}, since the YES bounds grow
linearly with the source order. For
total depth, start with a
$d$-regular instance $G$ on $n$ vertices from
Corollary~\ref{cor:regular-hardness}.
The identity $dn=2|E(G)|$ and oddness of $d$ imply that $n$ is even.
Add a perfect matching on the vertex set to make every degree even,
using auxiliary matching edges that may be parallel to original edges.
Each connected component of the augmented multigraph has even
degrees and therefore has an Euler tour, a closed walk traversing
every edge once. Orient its edges along this tour, then delete the
matching edges. Each vertex has outdegree
$(d+1)/2$ before deletion, so its outdegree in the original graph is
at least $(d-1)/2$.

Assign the internal vertices of each replacement path to the tail
of its oriented edge. For each vertex $v$, let $X_v$ consist of $v$
and the internal vertices assigned to it. The sets $X_v$ partition
$V(G_\ell)$. Each is connected, and
\begin{equation}\label{eq:subdivision-block-size}
 s=1+(\ell-1)(d-1)/2,\qquad |X_v|\ge s\quad(v\in V(G)).
\end{equation}
Each oriented edge $uv$ has a replacement path whose last edge
joins $X_u$ to $X_v$. Lemma~\ref{lem:connected-part-projection}
therefore projects any elimination forest $F$ of $G_\ell$ to an
elimination forest $\overline F$ of $G$, and gives
\begin{equation}\label{eq:subdivision-depth-projection}
 \totd(F)
 \overset{\text{Lem.~\ref{lem:connected-part-projection}}}{\ge}
 \sum_{v\in V(G)}|X_v|\operatorname{depth}_{\overline F}(v)
 \overset{\eqref{eq:subdivision-block-size}}{\ge}
 s\totd(\overline F)\ge s\totd(G).
\end{equation}

For the upper bound, use the same midpoint-recursive trees below the
deeper endpoints in an optimal elimination forest $F$ of $G$, and call
the resulting forest of $G_\ell$ $\widehat F$. Each of the
$\ell-1$ new vertices $x$ on a replacement path $uv$ satisfies
\begin{equation}\label{eq:subdivision-new-depth}
 \operatorname{depth}_{\widehat F}(x)\le
 \operatorname{depth}_F(\text{deeper endpoint of }uv)+h.
\end{equation}
Each vertex is the deeper endpoint of at
most $d$ edges, so summing the original and new depths gives
\[
\begin{aligned}
 s\totd(G)
 &\overset{\eqref{eq:subdivision-depth-projection}}{\leq}
 \totd(G_\ell)\\
 &\overset{\eqref{eq:subdivision-new-depth}}{\leq}
 \totd(G)+(\ell-1)\sum_{uv\in E(G)}
       \operatorname{depth}_{F}(\text{deeper endpoint of }uv)
       +h(\ell-1)|E(G)|\\
 &\leq
 [1+d(\ell-1)]\totd(G)+h(\ell-1)|E(G)|.
\end{aligned}
\]
Since $d\geq3$, using the definition of $s$, $|E(G)|=dn/2$,
and $C=h(\ell-1)d/2$ gives
\[
 3s-[1+d(\ell-1)]=2+\frac{(d-3)(\ell-1)}2\geq0,\qquad
 s\totd(G)\leq\totd(G_\ell)\leq3s\totd(G)+Cn.
\]
Choose a regular SSE source with a gap larger than three times
the target factor. This fixes $d$; next fix $\ell$, and then enlarge
the source so that $Cn$ is negligible compared with its quadratic
YES scale. The resulting gap survives by the calculation in
Appendix~\ref{app:transfer-subdivision}.

\end{proof}

Equation~\eqref{eq:tp-adjusted-cost} also transfers the total-depth
hardness conclusions of both corollaries to trivially perfect
completion.

